\section{Παραδοσιακές Προσεγγίσεις στην Ανάλυση Αρχείων Καταγραφής}

\subsection{Εισαγωγή}

Η ανάλυση αρχείων καταγραφής αποτελεί διαχρονικά βασικό πυλώνα της ασφάλειας δικτύων, επιτρέποντας την ανίχνευση ανωμαλιών, εισβολών και σφαλμάτων συστήματος. Πριν την εμφάνιση των σύγχρονων τεχνικών βαθιάς μάθησης, οι παραδοσιακές προσεγγίσεις βασίζονταν σε συστήματα βασισμένα σε κανόνες, ανίχνευση βάσει υπογραφών και κλασικές μεθόδους μηχανικής μάθησης. Οι τεχνικές αυτές έχουν υιοθετηθεί ευρέως σε συστήματα \en{Security Information and Event Management (SIEM)} και \en{intrusion detection systems (IDS)}, αποτελώντας τη βάση πολλών υπαρχουσών λύσεων κυβερνοασφάλειας.

Παρά την αποτελεσματικότητά τους σε ορισμένα σενάρια, οι παραδοσιακές μέθοδοι παρουσιάζουν σημαντικούς περιορισμούς όταν εφαρμόζονται σε σύγχρονα περιβάλλοντα \en{5G} και \en{6G}. Ο αυξανόμενος όγκος, η ταχύτητα και η πολυπλοκότητα των δεδομένων απαιτούν πιο επεκτάσιμες και προσαρμοστικές λύσεις. Το κεφάλαιο αυτό εξετάζει τις βασικές παραδοσιακές προσεγγίσεις και τους περιορισμούς τους.

\subsection{Συστήματα Βασισμένα σε Κανόνες}

Τα συστήματα βασισμένα σε κανόνες αποτελούν μία από τις πρώτες και πιο διαδεδομένες προσεγγίσεις για την ανάλυση αρχείων καταγραφής. Βασίζονται σε προκαθορισμένους κανόνες που περιγράφουν γνωστά πρότυπα κακόβουλης ή ανώμαλης συμπεριφοράς. Οι αναλυτές ασφάλειας ορίζουν αυτούς τους κανόνες χειροκίνητα. Τα συστήματα \en{SIEM} αποτελούν χαρακτηριστικό παράδειγμα. Συλλέγουν και συσχετίζουν δεδομένα καταγραφής από πολλαπλές πηγές και εφαρμόζουν κανόνες συσχέτισης για την ανίχνευση συμβάντων ασφάλειας \cite{behl2017cybersecurity}. Παρότι είναι εύκολα στην υλοποίηση και ερμηνεύσιμα, εξαρτώνται από χειροκίνητη συντήρηση και αδυνατούν να ανιχνεύσουν νέες ή εξελισσόμενες απειλές.

\subsection{Ανίχνευση Βάσει Υπογραφών}

Η ανίχνευση βάσει υπογραφών χρησιμοποιείται ευρέως σε συστήματα ανίχνευσης εισβολών και λογισμικά προστασίας από ιούς. Αναγνωρίζει απειλές συγκρίνοντας μοτίβα με μια βάση γνωστών υπογραφών \cite{garcia2009anomaly}. Το βασικό πλεονέκτημα είναι η υψηλή ακρίβεια για γνωστές απειλές με χαμηλές ψευδώς θετικές προβλεψεις. Ωστόσο, είναι αναποτελεσματική έναντι επιθέσεων μηδενικής ημέρας (\en{zero-day}) και άγνωστων απειλών. Σε σύγχρονα δίκτυα, όπου οι επιθέσεις εξελίσσονται συνεχώς, απαιτούνται συχνές ενημερώσεις των βάσεων υπογραφών.

\subsection{Προσεγγίσεις Μηχανικής Μάθησης}

Για την αντιμετώπιση των περιορισμών των προηγούμενων μεθόδων, έχουν αναπτυχθεί τεχνικές μηχανικής μάθησης (\en{ML}).

\subsubsection{Επιβλεπόμενη Μάθηση}

Οι μέθοδοι επιβλεπόμενης μάθησης χρησιμοποιούν επισημασμένα δεδομένα για την εκπαίδευση μοντέλων που ταξινομούν συμβάντα ως κανονικά ή κακόβουλα. Συνήθεις αλγόριθμοι περιλαμβάνουν δέντρα αποφάσεων, \en{support vector machines (SVM)} και νευρωνικά δίκτυα \cite{nguyen2018deep}. Παρουσιάζουν υψηλή ακρίβεια όταν υπάρχουν επαρκή δεδομένα, αλλά η συλλογή επισημασμένων δεδομένων είναι δύσκολη και δεν γενικεύονται εύκολα σε νέες επιθέσεις.

\subsubsection{Μη Επιβλεπόμενη Μάθηση}

Οι τεχνικές μη επιβλεπόμενης μάθησης δεν απαιτούν επισημασμένα δεδομένα. Εντοπίζουν πρότυπα και ανωμαλίες μέσω αλγορίθμων όπως το \en{k-means} \cite{garcia2009anomaly}. Είναι χρήσιμες για τον εντοπισμό άγνωστων επιθέσεων, αλλά συχνά παράγουν πολλά ψευδώς θετικά και απαιτούν επιπλέον ανάλυση.

\subsection{Περιορισμοί Παραδοσιακών Προσεγγίσεων}

Οι παραδοσιακές προσεγγίσεις παρουσιάζουν σημαντικούς περιορισμούς ως προς την εφαρμογή τους σε σύγχρονα δίκτυα μεγάλης κλίμακας. Αρχικά, εμφανίζουν περιορισμένη επεκτασιμότητα, καθώς δεν κλιμακώνονται εύκολα σε μεγάλα περιβάλλοντα και η επεξεργασία τεράστιου όγκου δεδομένων είναι ιδιαίτερα απαιτητική υπολογιστικά, ειδικά σε δίκτυα όπως το \en{5G}. Επιπλέον, βασίζονται σε μεγάλο βαθμό σε χειροκίνητη επιλογή χαρακτηριστικών (\en{feature engineering}), διαδικασία που απαιτεί εξειδικευμένη γνώση, σημαντικό χρόνο και συνεχή προσαρμογή \cite{he2017drain}. Τέλος, διαθέτουν περιορισμένη ικανότητα γενίκευσης, καθώς δυσκολεύονται να προσαρμοστούν σε νέες ή άγνωστες επιθέσεις, λειτουργώντας κυρίως αντιδραστικά και χωρίς να προσφέρουν επαρκή προστασία απέναντι σε προηγμένες απειλές.